\documentclass{homework}
\input{preamble}
\usepackage{hyperref}
\author{Jim Fowler}
\course{Math 5520H}
\title{Bessel functions}
\date{October 31, 2019}
\newcommand{\C}{\mathbb{C}}
\DeclareMathOperator{\image}{im}
\begin{document}
\maketitle


A Bessel functioni, of order $n$, is a solution to
\[
\displaystyle x^{2}{\frac {d^{2}y}{dx^{2}}}+x{\frac {dy}{dx}}+\left(x^{2}-n ^{2}\right)y=0.
\]

\begin{problem}
Show that $J_n(x)$, the Bessel function of the first kind of order $n$, has a series expansion
\[
  J_n(x) = \sum_{k=0}^\infty \frac{(-1)^k \, (x/2)^{2k + n}}{k! \, \Gamma(n + k + 1)}.
\]
\end{problem}

\begin{problem}
  Suppose $x \neq 0$ is a zero of $J_n$.  Then $x$ is a simple zero.
\end{problem}

\begin{problem}
  If $J_n(x) = 0$, then $J_n(-x) = 0$.
\end{problem}

\begin{problem}
  If $J_n(x) = 0$, then $J_n(\bar{x}) = 0$.
\end{problem}

\begin{problem}
  Differentiate
  \[
    x \left( b J_n(ax) J'_n(bx) - a J'_n(ax) J_n(bx) \right)
  \]
  to compute
   $(a^2-b^2) \int_0^x t J_n(at) J_n(bt) \, dt$.
\end{problem}

\vfill

\begin{problem}
  Apply the previous computation to the case that  $a$ and $b$ are conjugate zeros of $J_n(x)$.
\end{problem}

\vfill

\begin{problem}
  What can you conclude about the zeroes of $J_n$?
\end{problem}

\vfill

\end{document}
